% Generated draft: requires mathematical and source review.
\documentclass[11pt]{article}
\usepackage[margin=25mm]{geometry}
\usepackage{fontspec}
\setmainfont{Latin Modern Roman}
\usepackage{amsmath,amssymb,mathtools,unicode-math}
\setmathfont{Latin Modern Math}
\usepackage{enumitem,needspace,xurl,hyperref}
\setlength{\parindent}{0pt}
\setlength{\parskip}{5pt}
\setlength{\emergencystretch}{3em}
\newcommand{\N}{\mathbb{N}}
\newcommand{\Z}{\mathbb{Z}}
\newcommand{\Q}{\mathbb{Q}}
\newcommand{\R}{\mathbb{R}}
\newcommand{\C}{\mathbb{C}}
\newcommand{\F}{\mathbb{F}}
\newcommand{\A}{\mathbb{A}}
\newcommand{\PP}{\mathbb{P}}
\newcommand{\E}{\mathbb{E}}
\newcommand{\eps}{\varepsilon}
\newcommand{\dd}{\,\mathrm{d}}
\begin{document}
\section*{UnsolvedMath ID: 30004276}
\textbf{Source problem code:} catalogue entry
\subsection{Definitions and mathematical statement}
Let \(n\) denote the number of terms in the self-similar representation of a function \(P\). Assume, in the terminology of the supplied source, that \(P\) has zero spectral order and that its generalized derivative \(P'\) is a noncompact Sobolev multiplier. Consider the singular Sturm--Liouville boundary-value problem
\[
-y''=\lambda P'y,\qquad y(0)=y(1)=0.
\]
Here \(\lambda\) is the spectral parameter. The precise function spaces, meaning of the multiplier action \(P'y\), and operator or operator-pencil realization are not specified in the supplied record.

The requested problem is to determine the complete spectrum of this boundary-value problem, including its spectral type and eigenvalue behavior. More specifically, the supplied formulation asks whether, for every \(P\) satisfying the stated hypotheses, the spectrum consists of \(n-1\) possibly overlapping intervals of continuous spectrum together with eigenvalues, with at most one eigenvalue in each gap and at most \(n-2\) eigenvalues in total.

This is a conjectural description, not an asserted spectral theorem. The supplied record does not define an \(n\)-term self-similar function, zero spectral order, or the Sobolev spaces involved in the noncompact multiplier hypothesis. It also leaves unspecified the admissible self-similarity parameters, any reality or sign assumptions, the spectral realization, interval endpoint conventions, the meaning of a gap, and whether eigenvalues are counted with multiplicity. Consequently, the displayed problem and proposed spectral description retain the source terminology and require these qualifications to be resolved before becoming a fully specified mathematical assertion. The allowed range of \(n\) also requires clarification, since the proposed upper bound \(n-2\) is negative when \(n=1\).
\subsection{Short English statement}
For a self-similar weight whose generalized derivative is a noncompact Sobolev multiplier, determine the spectrum of the stated Dirichlet problem. The supplied conjecture predicts \(n-1\) possibly overlapping continuous-spectrum intervals, at most one eigenvalue in each gap, and at most \(n-2\) eigenvalues altogether. Essential definitions and realization details remain unresolved; this is a statement-only draft requiring human mathematical review.
\subsection{Sources}
\begin{itemize}
\item Local UnsolvedMath registry: supplied record with permanent ID \(30004276\), title ``Spectrum of Boundary Problems with Noncompact Multiplier Weights,'' exact source problem code \texttt{catalogue entry}, in the Oberwolfach Reports Open Problems collection. No individual proposer is supplied; the proposed year is recorded as 2019.
\item Supplied primary-source citation: \emph{Mini-Workshop: Self-adjoint Extensions in New Settings (2020)}, Oberwolfach Reports, pp.\ 2911--2950, DOI \texttt{10.4171/owr/2019/47}. \url{https://doi.org/10.4171/owr/2019/47}. The supplied citation uses 2020, whereas the DOI and proposed-year metadata contain 2019; their relationship has not been independently checked.
\item Supplied literature source: ``Sharov--Sheipak, string equation with a noncompact multiplier.'' \url{https://www.mathnet.ru/php/archive.phtml?issue=4&jrnid=aa&option_lang=eng&volume=33&wshow=issue&year=2021}. The supplied link identifies a journal issue rather than a complete article citation.
\item Supplied literature source: ``original Oberwolfach report.'' \url{https://doi.org/10.4171/owr/2019/47}.
\item Literature assessment dated exactly \texttt{2026-08-22}: the supplied curation labels the problem open and reports that later work determines continuous spectrum and eigenvalue behavior for important noncompact-multiplier string models, while the general complete spectral description remains unestablished. This is supplied dated curation, not an independently verified literature assessment; its time-sensitive status and cited claims require independent verification.
\item The supplied registry labels the statement \texttt{corrected\_verified} and reports that the official report was checked to replace a dangling boundary-value reference with the displayed equation, boundary conditions, self-similar weight hypothesis, and spectral conjecture. These are supplied verification claims, not verification performed for this draft.
\end{itemize}
\end{document}
